\documentclass[a4paper,12pt]{article}
\usepackage[utf8]{inputenc}
\usepackage{CJKutf8}
\usepackage[margin=2.2cm]{geometry}
\usepackage{array}
\usepackage{booktabs}
\usepackage{longtable}
\usepackage{tabularx}
\usepackage{enumitem}
\usepackage{hyperref}
\usepackage{xcolor}

\hypersetup{
  colorlinks=true,
  linkcolor=black,
  urlcolor=blue
}

\setlist[itemize]{leftmargin=1.8em}
\setlist[enumerate]{leftmargin=2.0em}
\renewcommand{\arraystretch}{1.2}

\title{基于 ESP32 的多 Agent 智能管家系统功能说明}
\author{}
\date{}

\begin{document}
\begin{CJK*}{UTF8}{gbsn}
\maketitle

\begin{abstract}
本文基于项目当前仓库实现，对多 Agent 智能管家系统的功能进行系统化梳理，并按作品说明文档格式输出为 \LaTeX{} 版本。系统以 ESP32-S3 为核心，构建了 coordinator\_agent、sensor\_agent、control\_agent、guardian\_agent 四角色协同架构，结合 ESP32-P4/网页前端展示终端，形成“自然语言入口---多节点调度---环境感知---硬件执行---策略审计---结果回传”的完整闭环。系统已实现 Feishu/WebSocket/串口多入口交互、MQTT Mesh 节点协同、环境感知、执行器控制、自动化工作流、条件规则、定时任务、技能装载、长期记忆、子代理调用、可视化时间线、Sandbox 拦截展示、Guardian 策略裁决和控制急停等能力。本文重点列出当前代码中已经落地的功能边界、工具能力、软件组成与扩展方向，便于后续参赛文档、答辩汇报与系统归档统一使用。
\end{abstract}

\section{作品概述}

\subsection{功能与特性}
本系统的功能可概括为以下九类：
\begin{enumerate}
  \item \textbf{多入口自然语言交互}：支持 Feishu 机器人、网页 WebSocket 控制台以及串口 CLI 作为用户入口。
  \item \textbf{四角色 Agent Mesh 协同}：Coordinator 负责任务理解与调度，Sensor 负责环境采集，Control 负责硬件执行，Guardian 负责策略裁决与审计。
  \item \textbf{环境感知}：支持温湿度、空气质量、光照、人体存在、超声距离等数据采集，并支持运行时虚拟设备读取。
  \item \textbf{执行器控制}：支持 WS2812 状态灯、普通 GPIO、GPIO4 加湿器、GPIO5 风扇、GPIO6 独立 LED、舵机、红外空调和音频蜂鸣测试。
  \item \textbf{自动化执行}：支持多步工作流、条件规则和定时任务三种自动化机制。
  \item \textbf{安全控制}：支持 Sandbox、参数白名单、Guardian policy gate、控制急停与状态回读。
  \item \textbf{智能运行时能力}：支持 Skills、长期记忆、用户画像、子代理、Lua 脚本扩展和运行时设备清单。
  \item \textbf{可视化展示}：支持节点拓扑、时间线、Sandbox 事件流、Skills Studio、用户偏好面板与 AI 通信面板。
  \item \textbf{运维与扩展}：支持 SPIFFS 文件管理、运行时 Skills 安装和串口调试命令。
\end{enumerate}

\subsection{应用领域}
本系统面向以下应用场景：
\begin{itemize}
  \item \textbf{智能家居控制}：通过自然语言完成风扇、加湿器、灯光、空调等联动控制。
  \item \textbf{实验室与宿舍环境管理}：持续监测温湿度、CO\textsubscript{2}、TVOC、光照和人体靠近状态，并触发规则动作。
  \item \textbf{边缘智能教学平台}：展示从 LLM 到工具调用、从 Mesh 调度到硬件执行的完整闭环。
  \item \textbf{多 Agent 协同研究}：验证多角色分工、策略裁决、时间线审计与运行时技能加载机制。
  \item \textbf{AIoT 原型验证平台}：快速接入 I2C/UART/SPI/GPIO 设备，验证新传感器与新执行器能力。
\end{itemize}

\subsection{主要技术特点}
\begin{itemize}
  \item 基于 ESP-IDF、FreeRTOS 和 cJSON 构建轻量级 Agent Runtime。
  \item 基于 MQTT Mesh 实现跨节点消息路由、命令下发、结果回传和时间线广播。
  \item 基于 ReAct 风格工具调用，实现从自然语言到结构化工具执行的转换。
  \item 基于 SPIFFS 实现 Skills、Memory、Profile、脚本和设备清单的本地持久化。
  \item 基于 Guardian + Sandbox + 白名单参数校验实现安全执行边界。
  \item 基于前端聚合服务，实现 MQTT 数据聚合、技能安装代理和网页聊天网关。
  \item 基于自动化引擎实现工作流、条件规则和定时任务的后台运行。
\end{itemize}

\subsection{主要性能指标}
\begin{table}[h]
\centering
\begin{tabular}{p{4.2cm}p{8.6cm}}
\toprule
\textbf{指标项} & \textbf{当前实现} \\
\midrule
Agent 角色数 & 4 个固定 ESP32-S3 角色节点：coordinator\_agent、sensor\_agent、control\_agent、guardian\_agent \\
展示终端 & 1 个显示/前端终端，支持网页控制台与状态展示 \\
注册工具能力 & 48 项已注册工具能力（以 \texttt{main/tools/tool\_registry.c} 为准） \\
Slash 指令数 & 22 项显式 Slash 命令/动作指令 \\
环境感知能力 & 温湿度、eCO\textsubscript{2}、TVOC、光照、人体存在、超声距离、虚拟只读设备 \\
执行器能力 & WS2812、普通 GPIO、加湿器、风扇、独立 LED、舵机、Gree 空调 IR、MAX98357 音频、虚拟控制设备 \\
自动化能力 & 工作流 1 类、条件规则 1 类、定时任务 3 类（every/at/daily） \\
持久化能力 & Skills、Memory、Profile、脚本、设备清单、会话摘要均可落盘到 SPIFFS \\
用户入口 & Feishu、WebSocket 网页前端、串口 CLI \\
内部协同链路 & MQTT Mesh + Timeline + StateBoard + PolicyDecision \\
\bottomrule
\end{tabular}
\end{table}

\subsection{主要创新点}
\begin{enumerate}
  \item 将四角色 Agent 分工、MQTT Mesh 协作与硬件控制闭环统一在 ESP32 侧完成。
  \item 通过 Guardian、Sandbox 和工具白名单，把自然语言控制约束在可审计、可回退的安全边界内。
  \item 支持运行时 Skills 装载，无需重新烧录即可扩展板端行为知识。
  \item 将多步时序动作、条件规则和定时任务统一纳入后台自动化引擎，降低 LLM 持续占用。
  \item 在前端同时展示节点拓扑、时间线、Sandbox 事件和 AI 通信，使推理链路可视化。
\end{enumerate}

\subsection{设计流程}
系统设计流程如下：首先根据任务需求将系统划分为调度、感知、执行、安全四类职责；然后围绕 ESP32-S3 建立统一 Tool Registry，把各类传感器、执行器与系统能力封装为工具；随后通过 MQTT Mesh 连接多节点，并接入 Feishu、WebSocket 和串口入口；在此基础上补充 Skills、Memory、自动化引擎、Guardian 与前端展示；最后通过串口脚本、前端时间线和实际硬件动作进行联调验证。该流程保证了系统从单节点能力到多节点协作、再到可视化与运维能力的逐层演进。

\section{系统组成及功能说明}

\subsection{整体介绍}
系统总体架构可分为六层：
\begin{enumerate}
  \item \textbf{用户入口层}：Feishu、网页前端、串口 CLI；
  \item \textbf{Coordinator 调度层}：负责 LLM 对话、工具调用与 Mesh 调度；
  \item \textbf{Sensor 感知层}：负责温湿度、空气质量、光照、人体存在、距离与虚拟设备读取；
  \item \textbf{Control 执行层}：负责 GPIO、灯光、风扇、加湿器、空调、舵机和音频输出；
  \item \textbf{Guardian 安全层}：负责策略裁决、审计、状态板与 Watchdog；
  \item \textbf{展示与扩展层}：包括网页控制台、ESP32-P4 展示终端与运行时 Skills。
\end{enumerate}

\begin{table}[h]
\centering
\begin{tabular}{p{3.2cm}p{3.2cm}p{7.0cm}}
\toprule
\textbf{节点/模块} & \textbf{角色} & \textbf{主要功能} \\
\midrule
Coordinator 节点 & 调度中枢 & 接收用户消息、构建上下文、调用工具、下发 Mesh 命令并生成回复 \\
Sensor 节点 & 感知中枢 & 周期采集环境数据，响应传感器读取命令，发布 telemetry \\
Control 节点 & 执行中枢 & 执行灯光、GPIO、舵机、风扇、加湿器、空调和音频等控制动作 \\
Guardian 节点 & 安全中枢 & 对高风险操作进行 policy 决策、审计、StateBoard 更新和节点看门狗 \\
Display/Web 前端 & 展示终端 & 展示拓扑、时间线、Sandbox、Skills、偏好和聊天消息 \\
\bottomrule
\end{tabular}
\end{table}

\subsection{硬件系统介绍}

\subsubsection{硬件整体介绍}
硬件系统由多块 ESP32-S3 节点和一个展示终端组成。各节点通过 MQTT Mesh 协同工作，其中 coordinator 节点承载 Agent 主循环，sensor 节点承载环境传感器，control 节点承载执行器，guardian 节点承载安全边界和审计逻辑。系统外围可进一步连接 ESP32-P4/C6 显示终端。

\subsubsection{机械设计介绍}
当前系统核心侧重电子与软件联调，机械结构不是主要实现对象，因此暂未形成独立的机械机构设计模块。若后续参赛版本加入外壳、传感器支架或桌面终端结构，可在本节补充总体外观、局部支撑件和装配图。

\subsubsection{电路各模块介绍}
\begin{itemize}
  \item \textbf{温湿度模块}：AHT10/AHT20，通过 I2C 读取温度与湿度。
  \item \textbf{空气质量模块}：SGP30，读取 eCO\textsubscript{2} 与 TVOC。
  \item \textbf{光照模块}：GY-30/BH1750，读取环境照度。
  \item \textbf{人体/距离模块}：支持数字人体传感器输出，也支持 HC-SR05 超声测距。
  \item \textbf{灯光模块}：WS2812 状态灯，支持命名颜色和 RGB 值控制。
  \item \textbf{GPIO 负载模块}：GPIO4 对应加湿器，GPIO5 对应风扇，GPIO6 对应独立 LED。
  \item \textbf{执行机构模块}：舵机 PWM 控制、Gree 空调红外发射。
  \item \textbf{音频模块}：MAX98357 I2S 扬声器测试输出。
  \item \textbf{扩展设备模块}：支持运行时设备清单驱动的 UART/I2C/SPI/ADC/GPIO 外设接入。
\end{itemize}

\subsection{软件系统介绍}

\subsubsection{软件整体介绍}
软件系统由板端 Agent Runtime、Mesh 通信层、自动化引擎、Skills/Memory 持久层、前端聚合服务和网页控制台共同组成。板端负责实时感知、执行和策略判断；前端负责聚合 MQTT 数据、代理聊天网关、展示时间线和管理运行时 Skills。整体软件形成了“入口---理解---执行---审计---展示”的端到端闭环。

\subsubsection{软件各模块介绍}

\paragraph{1) Agent Runtime 与上下文构建}
\begin{itemize}
  \item 支持 LLM 主循环、上下文拼接、工具调用、结果归纳和最终回复。
  \item 支持 Tool Registry、Context Builder、Session/Brief/Trace 管理。
  \item 支持长期 Memory、用户 Profile 和记忆摘要注入。
\end{itemize}

\paragraph{2) 多节点协同与消息路由}
\begin{itemize}
  \item 支持 MQTT Mesh 命令、结果、timeline、state、telemetry、stateboard。
  \item 支持按 \texttt{target\_role} 或 \texttt{target\_node} 定向路由。
  \item 支持 control\_state、急停、解急停、agent\_task 等远程动作。
\end{itemize}

\paragraph{3) 自动化与定时调度}
\begin{itemize}
  \item \textbf{Workflow}：执行带时延的多步动作序列。
  \item \textbf{Rule}：根据温度/湿度阈值持续触发条件动作。
  \item \textbf{Cron}：支持周期、定点和每日定时注入任务。
\end{itemize}

\paragraph{4) Skills、Lua 与运行时扩展}
\begin{itemize}
  \item Skills 从 \texttt{/spiffs/skills/} 加载，可列出、查看、安装和摘要注入。
  \item 支持 Skill 观察记录、运行时 Skill 安装与前端 Skills Studio。
  \item 支持 Lua 运行时信息查询、脚本列出、同步/异步运行、任务查询与停止。
  \item 支持运行时设备清单，统一描述简单传感器与执行器协议。
\end{itemize}

\paragraph{5) 前端与展示层}
\begin{itemize}
  \item 支持节点总览、能力目录、环境卡片、调度时间线、Sandbox 面板、用户偏好面板。
  \item 支持 AI 通信窗口、聊天网关状态显示和运行时 Skill 管理。
  \item 支持 admin/operator/viewer 三类演示角色。
\end{itemize}

\paragraph{6) 安全与审计层}
\begin{itemize}
  \item 支持 Sandbox 拦截，敏感写入和持久规则要求显式确认。
  \item 支持 Guardian policy\_decision、StateBoard、Watchdog 和风险事件追踪。
  \item 支持 control\_agent 急停闭锁，防止高频或异常动作持续执行。
\end{itemize}

\paragraph{7) 48 项工具能力清单}
\begin{longtable}{p{2.8cm}p{3.6cm}p{8.0cm}}
\toprule
\textbf{类别} & \textbf{工具名} & \textbf{功能说明} \\
\midrule
\endfirsthead
\toprule
\textbf{类别} & \textbf{工具名} & \textbf{功能说明} \\
\midrule
\endhead
信息服务 & \texttt{web\_search} & 联网搜索当前信息 \\
信息服务 & \texttt{get\_weather} & 查询当前位置或指定地点天气 \\
信息服务 & \texttt{get\_current\_time} & 获取并同步系统时间 \\
Lua 扩展 & \texttt{lua\_runtime\_info} & 查询 Lua 运行时可用性和限制 \\
Lua 扩展 & \texttt{lua\_list\_modules} & 列出 Lua 可用模块 \\
Lua 扩展 & \texttt{lua\_list\_scripts} & 列出可运行脚本 \\
Lua 扩展 & \texttt{lua\_run\_source} & 运行内联 Lua 源码 \\
Lua 扩展 & \texttt{lua\_run\_script} & 运行 SPIFFS 中的 Lua 脚本 \\
Lua 扩展 & \texttt{lua\_run\_script\_async} & 后台异步运行 Lua 脚本 \\
Lua 扩展 & \texttt{lua\_list\_jobs} & 列出 Lua 后台任务 \\
Lua 扩展 & \texttt{lua\_get\_job} & 查询单个 Lua 后台任务 \\
Lua 扩展 & \texttt{lua\_stop\_job} & 停止 Lua 后台任务 \\
记忆与技能 & \texttt{memory\_profile\_set} & 写入稳定用户画像和长期偏好 \\
记忆与技能 & \texttt{skill\_observation\_add} & 记录技能验证结果与观察 \\
感知 & \texttt{read\_temperature\_humidity} & 读取本地或远程温湿度 \\
协同 & \texttt{mesh\_send\_command} & 发布标准 MQTT Mesh 命令 \\
扩展设备 & \texttt{virtual\_device\_read} & 读取运行时设备清单中的只读设备 \\
扩展设备 & \texttt{virtual\_device\_control} & 控制运行时设备清单中的可控设备 \\
自动化 & \texttt{automation\_create\_workflow} & 创建并启动多步工作流 \\
自动化 & \texttt{automation\_create\_rule} & 创建持续条件规则 \\
自动化 & \texttt{automation\_list} & 列出工作流和规则 \\
自动化 & \texttt{automation\_remove} & 删除工作流或规则 \\
Agent 协作 & \texttt{spawn\_subagent} & 派生临时子代理处理聚焦子任务 \\
综合感知 & \texttt{read\_environment} & 一次性读取温湿度、空气质量和光照 \\
综合感知 & \texttt{read\_presence} & 读取人体存在或近距靠近状态 \\
综合感知 & \texttt{hc\_sr05\_read\_distance} & 直接读取超声距离 \\
文件系统 & \texttt{read\_file} & 读取 SPIFFS 文件 \\
文件系统 & \texttt{write\_file} & 写入 SPIFFS 文件 \\
文件系统 & \texttt{edit\_file} & 修改 SPIFFS 文件内容 \\
文件系统 & \texttt{list\_dir} & 列出 SPIFFS 目录内容 \\
基础控制 & \texttt{gpio\_write} & 控制普通 GPIO 电平 \\
基础控制 & \texttt{copper\_gpio\_write} & 控制 GPIO4/5/6 专用铜皮引脚 \\
基础控制 & \texttt{set\_humidifier} & 控制加湿器开关 \\
基础控制 & \texttt{set\_fan} & 控制风扇开关 \\
基础控制 & \texttt{set\_device\_led} & 控制 GPIO6 独立 LED \\
基础控制 & \texttt{gpio\_read} & 读取单个 GPIO 电平 \\
基础控制 & \texttt{gpio\_read\_all} & 读取全部允许 GPIO 电平 \\
灯光控制 & \texttt{ws2812\_set} & 控制 WS2812 RGB 数值 \\
灯光控制 & \texttt{set\_status\_light} & 以颜色名控制状态灯 \\
执行机构 & \texttt{servo\_write} & 控制舵机角度或脉宽 \\
执行机构 & \texttt{gree\_ac\_control} & 控制格力空调红外命令 \\
音频输出 & \texttt{max98357\_play\_tone} & 播放短音调测试音频链路 \\
空气质量 & \texttt{sgp30\_read\_air\_quality} & 直接读取 SGP30 eCO\textsubscript{2}/TVOC \\
空气质量 & \texttt{read\_air\_quality} & 高层空气质量读取接口 \\
光照感知 & \texttt{read\_light\_level} & 读取光照强度 \\
定时任务 & \texttt{cron\_add} & 新增周期/定点/每日定时任务 \\
定时任务 & \texttt{cron\_list} & 列出当前定时任务 \\
定时任务 & \texttt{cron\_remove} & 删除指定定时任务 \\
\bottomrule
\end{longtable}

\paragraph{8) Slash 指令能力}
系统还支持显式 Slash 路由与动作指令，包括：\texttt{/help}、\texttt{/clear}、\texttt{/clear\_all\_memory}、\texttt{/context\_status}、\texttt{/skills\_list}、\texttt{/skills\_show}、\texttt{/sensor}、\texttt{/control}、\texttt{/guardian}、\texttt{/subagent}、\texttt{/workflow}、\texttt{/rule}、\texttt{/local}、\texttt{/mesh}、\texttt{/status}、\texttt{/stop}、\texttt{/resume}、\texttt{/device}、\texttt{/profile}、\texttt{/skills}、\texttt{/privacy}、\texttt{/lua} 和 \texttt{/trace}。这些指令可直接约束 Agent 的路由、执行范围和工具选择。

\section{完成情况及性能参数}

\subsection{整体介绍}
从当前实现看，系统已经具备完整的多 Agent 智能家居闭环：用户可从 Feishu、网页或串口发起请求；Coordinator 负责理解任务并调用工具或下发 Mesh 命令；Sensor 节点持续上报环境数据；Control 节点执行灯光、风扇、加湿器、舵机等动作；Guardian 节点完成策略审计；前端统一展示拓扑、时间线、Sandbox 和 AI 通信结果。整体系统已具备可演示、可扩展、可归档的工程形态。

\subsection{工程成果}

\subsubsection{机械成果}
当前版本没有将机械结构作为主要成果，故该部分以“预留”处理。若后续引入壳体、面板、桌面终端支架或传感器固定结构，可在此补充实物图与装配图。

\subsubsection{电路成果}
当前电路成果主要包括：
\begin{itemize}
  \item AHT20 温湿度采集链路；
  \item SGP30 空气质量采集链路；
  \item BH1750/GY-30 光照采集链路；
  \item HC-SR05 超声/人体存在采集链路；
  \item WS2812 状态灯控制链路；
  \item GPIO4/5/6 对应加湿器、风扇和独立 LED 控制链路；
  \item 舵机 PWM 控制链路；
  \item Gree 空调红外发送链路；
  \item MAX98357 音频播放测试链路。
\end{itemize}

\subsubsection{软件成果}
当前软件成果主要包括：
\begin{itemize}
  \item ESP32 侧 LLM Agent Runtime；
  \item 四角色 MQTT Mesh 协同框架；
  \item 自动化工作流、规则、定时任务引擎；
  \item Skills Loader、Memory/Profile 持久化模块；
  \item Guardian Sandbox、StateBoard、Watchdog 审计链路；
  \item 前端控制台、聊天网关、技能安装代理；
\end{itemize}

\subsection{特性成果}
本系统当前已经可以量化展示的能力成果如下：
\begin{enumerate}
  \item 支持 4 个固定角色节点协同运行，形成分工明确的边缘多 Agent 架构。
  \item 支持 48 项注册工具能力，覆盖感知、控制、自动化、技能、文件和运行时扩展能力。
  \item 支持 3 类自动化机制：多步工作流、条件规则、定时任务。
  \item 支持 3 类用户入口：Feishu、WebSocket 网页前端、串口 CLI。
  \item 支持 22 项 Slash 指令，实现显式路由、显式清理、显式状态查询与显式技能访问。
  \item 支持 Runtime Skill 安装到 SPIFFS，无需重复烧录即可扩展系统知识。
  \item 支持 Sandbox 可视化，前端可以单独查看拦截事件、风险事件和确认要求。
  \item 支持控制急停与解除急停，适合在执行器联调阶段进行安全兜底。
\end{enumerate}

\section{总结}

\subsection{可扩展之处}
系统后续可沿以下方向扩展：
\begin{itemize}
  \item 扩展更多运行时设备清单，使 UART/I2C/SPI 外设无需重写大量固件即可接入；
  \item 继续优化前端交互链路和板端任务调度策略，降低实时控制与长上下文推理之间的资源竞争；
  \item 增加更丰富的用户画像与偏好闭环，实现更主动的环境控制与消息提醒。
\end{itemize}

\subsection{心得体会}
本系统的核心经验在于：面向真实硬件的 Agent 系统，不能只追求“能回答”，更要追求“能执行、能回看、能审计、能停下”。因此在设计中引入了角色分工、Mesh 协同、自动化引擎、Guardian 审核、Sandbox 拦截和前端时间线展示等机制，使自然语言控制不再停留在单轮对话，而是能够落到实际设备、后台任务和持续规则中。同时，系统保留了 Skills、Lua 和设备清单等扩展边界，为后续从单一智能家居演示走向复杂边缘智能平台留下了接口。

\section{参考文献}
\begin{enumerate}
  \item 项目文档：\texttt{docs/RESUME\_PROJECT\_INTRO.md}.
  \item 项目文档：\texttt{docs/PROJECT\_STRUCTURE.md}.
  \item 项目文档：\texttt{docs/LLM\_TO\_HARDWARE\_RUNTIME.md}.
  \item 前端文档：\texttt{frontend/README.md}.
  \item Espressif Systems. \textit{ESP-IDF Programming Guide}. \url{https://docs.espressif.com/projects/esp-idf/en/latest/esp32/}.
  \item MQTT Version 3.1.1 Specification. OASIS Standard.
  \item Feishu Open Platform Documentation. \url{https://open.feishu.cn/document/}.
  \item React Documentation. \url{https://react.dev/}.
  \item Ant Design Documentation. \url{https://ant.design/}.
\end{enumerate}

\end{CJK*}
\end{document}
